% =====================================================================
%  Lumina — BTP-2 Report
%  CS47006 (Project Part-II) | Spring 2025-2026 | IIT Kharagpur
%  Student : Abhinav Kumar Singh (22CS30005)
%  Supervisor : Professor Soumya Kanti Ghosh
% =====================================================================

\documentclass[12pt,a4paper]{report}

% ======================== PACKAGES ===================================
\usepackage[margin=1in]{geometry}
\usepackage[utf8]{inputenc}
\usepackage[T1]{fontenc}
\usepackage{lmodern}
\usepackage{times}
\usepackage{microtype}
\usepackage{setspace}
\usepackage{parskip}
\usepackage{ragged2e}
\usepackage{graphicx}
\usepackage{float}
\usepackage{caption}
\usepackage{booktabs}
\usepackage{array}
\usepackage{amsmath}
\usepackage{enumitem}
\usepackage{titlesec}
\usepackage{listings}
\usepackage{xcolor}
\usepackage{verbatim}
\usepackage[hidelinks]{hyperref}
\usepackage{url}

% ======================== CODE STYLE =================================
\lstset{
  language=Python,
  basicstyle=\ttfamily\small,
  breaklines=true,
  breakatwhitespace=true,
  frame=single,
  numbers=left,
  numberstyle=\tiny\color{gray},
  keywordstyle=\color{blue},
  commentstyle=\color{green!50!black},
  stringstyle=\color{orange}
}

% ======================== MISC SETUP =================================
\sloppy
\setlength{\parindent}{0pt}
\graphicspath{{figures/}}

% =====================================================================
\begin{document}

% =====================================================================
%  TITLE PAGE
% =====================================================================
\thispagestyle{empty}
\begin{center}
    \vspace*{0.8cm}
    \includegraphics[width=3.5cm]{kgplogo.png}\\[0.8cm]

    {\large \textbf{INDIAN INSTITUTE OF TECHNOLOGY, KHARAGPUR}}\\[0.3cm]
    {\normalsize \textbf{DEPARTMENT OF COMPUTER SCIENCE AND ENGINEERING}}\\[1.5cm]

    {\normalsize \textbf{Project-II (CS47006) Report submitted to}}\\[0.2cm]
    {\normalsize \textbf{Indian Institute of Technology Kharagpur}}\\[0.2cm]
    {\normalsize \textbf{in partial fulfilment for the award of the degree of}}\\[0.5cm]
    {\large \textbf{Bachelor of Technology and Master of Technology}}\\[0.2cm]
    {\normalsize \textbf{in}}\\[0.2cm]
    {\large \textbf{Computer Science and Engineering}}\\[1.5cm]

    \textbf{by}\\[0.3cm]
    {\large \textbf{Abhinav Kumar Singh}}\\[0.1cm]
    {\normalsize (22CS30005)}\\[1cm]

    \textbf{Under the supervision of}\\[0.3cm]
    {\large \textbf{Professor Soumya Kanti Ghosh}}\\[1.5cm]

    {\large \textbf{Lumina: Multi-Model Forecasting of Night-Time Light Data}}\\[1.5cm]

    \textbf{April 2026}
\end{center}

% =====================================================================
%  DECLARATION
% =====================================================================
\newpage
\thispagestyle{empty}
\begin{center}
    \hypertarget{sec:declaration}{}%
    {\Large \textbf{DECLARATION}}\\[1.5cm]
\end{center}

\onehalfspacing
I certify that:\\[0.3cm]
(a) The work contained in this report has been done by me under the guidance of my supervisor.\\[0.3cm]
(b) The work has not been submitted to any other Institute for any degree or diploma.\\[0.3cm]
(c) I have conformed to the norms and guidelines given in the Ethical Code of Conduct of the Institute.\\[0.3cm]
(d) Whenever I have used materials (data, theoretical analysis, figures, and text) from other sources, I have given due credit to them by citing them in the text of the thesis and giving their details in the references. Further, I have taken permission from the copyright owners of the sources, whenever necessary.\\[2cm]

\begin{flushright}
Date: April 2026\\
Place: Kharagpur\\[1.5cm]
\textbf{(Abhinav Kumar Singh)}\\
(22CS30005)
\end{flushright}

% =====================================================================
%  CERTIFICATE
% =====================================================================
\newpage
\thispagestyle{empty}
\begin{center}
    {\Large \textbf{DEPARTMENT OF COMPUTER SCIENCE AND ENGINEERING}}\\[0.5cm]
    {\large \textbf{INDIAN INSTITUTE OF TECHNOLOGY, KHARAGPUR}}\\
    KHARAGPUR -- 721302, INDIA\\[2cm]

    \hypertarget{sec:certificate}{}%
    {\Large \textbf{CERTIFICATE}}\\[1.5cm]
\end{center}

This is to certify that the project report entitled \textbf{``Lumina: Multi-Model Forecasting of Night-Time Light Data''}, submitted by \textbf{Abhinav Kumar Singh (Roll No.\ 22CS30005)} to the Indian Institute of Technology Kharagpur, towards partial fulfilment of the requirements for the award of the degree of \textbf{Bachelor of Technology and Master of Technology in Computer Science and Engineering}, is a record of bona fide work carried out by him under my supervision and guidance during the Spring Semester 2025--26.\\[1.5cm]

\begin{flushright}
Date: April 2026\\
Place: Kharagpur\\[1cm]
\textbf{Professor Soumya Kanti Ghosh}\\
Department of Computer Science and Engineering\\
Indian Institute of Technology Kharagpur\\
Kharagpur -- 721302, India
\end{flushright}

% =====================================================================
%  ABSTRACT
% =====================================================================
\newpage
\thispagestyle{empty}

{\centering
    \hypertarget{sec:abstract}{}%
    {\Large \textbf{ABSTRACT}}\\[0.5cm]
    \noindent\rule{\linewidth}{0.6pt}\\[1cm]
}

\textbf{Name of the Student:} Abhinav Kumar Singh\\
\textbf{Roll No:} 22CS30005\\
\textbf{Degree:} Bachelor of Technology and Master of Technology\\
\textbf{Department:} Department of Computer Science and Engineering\\
\textbf{Thesis Title:} Lumina: Multi-Model Forecasting of Night-Time Light Data\\
\textbf{Thesis Supervisor:} Professor Soumya Kanti Ghosh\\
\textbf{Month and Year of Submission:} April 2026\\[0.8cm]

\noindent\rule{\linewidth}{0.6pt}\\[1cm]

Nighttime light (NTL) satellite imagery provides a globally consistent, high-frequency proxy for
human activity, economic output, and urban development. This thesis presents \textbf{Lumina}, a
multi-model NTL forecasting platform developed as BTP-2 at IIT Kharagpur. The platform ingests
147 monthly VIIRS VCMSLCFG cloud-free composites (January 2014 -- March 2026) for two West Bengal
cities --- Kharagpur (35$\times$45 pixel grid) and Kolkata (70$\times$70 pixel grid) --- sourced via
the Google Earth Engine public data catalogue.

Three forecasting models of increasing spatial and architectural complexity are implemented and
evaluated. First, SARIMA serves as the classical statistical baseline, with orders selected via
AIC grid search: SARIMA(0,1,1)(0,1,1)[12] for Kharagpur (AIC = 189.95) and
SARIMA(0,1,2)(0,1,1)[12] for Kolkata (AIC = 444.22). Second, a stacked two-layer LSTM with Monte
Carlo Dropout furnishes probabilistic scalar forecasts with uncertainty intervals. Third, a
ConvLSTM encoder-decoder performs pixel-level spatial forecasting, preserving the full spatial
structure of the NTL field at each forecast step.

All models are validated against a 24-month held-out test set (January 2024 -- December 2025) and,
critically, against three months of real 2026 VIIRS observations (January -- March 2026) downloaded
in April 2026 --- a true live validation. On the live set, SARIMA achieves mean absolute percentage
errors of 10.0\% (Kharagpur) and 16.2\% (Kolkata), outperforming LSTM (17.5\%, 22.8\%) and
ConvLSTM (15.5\%, 29.9\%).

All results, raw satellite tiles, forecast outputs, and interactive visualisations are exposed
through the \textbf{Lumina} Streamlit dashboard (8 pages, publicly deployed at
\texttt{https://lumina-iamabhinav-dev.streamlit.app}). Future work proposes a multi-modal research paper
integrating NTL with WorldPop population estimates and district GDP proxies for pan-India
economic forecasting.

% =====================================================================
%  ACKNOWLEDGEMENTS
% =====================================================================
\newpage
\thispagestyle{empty}
\begin{center}
    \hypertarget{sec:ack}{}%
    {\Large \textbf{ACKNOWLEDGEMENTS}}\\[1cm]
\end{center}

I would like to express my sincere gratitude to my thesis supervisor, \textbf{Professor Soumya Kanti Ghosh}, Department of Computer Science and Engineering, IIT Kharagpur, for his continuous guidance, encouragement, and constructive feedback throughout both BTP-1 and BTP-2. His insights on remote sensing data and machine learning methods greatly shaped the direction of this project.

I am thankful to the faculty and staff of the Department of Computer Science and Engineering, IIT Kharagpur, for providing the academic environment and computational resources that made this work possible.

I am grateful to Google Earth Engine for providing free access to the VIIRS VCMSLCFG satellite imagery dataset, and to Streamlit Community Cloud for hosting the Lumina dashboard.

Finally, I express my deepest gratitude to my family --- my parents and my sister --- for their unwavering support and encouragement throughout my B.Tech--M.Tech programme. Their belief in my work has been a constant source of motivation.

\hfill \textbf{Abhinav Kumar Singh}\\
\hfill April 2026

% =====================================================================
%  TABLE OF CONTENTS (manual)
% =====================================================================
\newpage
\thispagestyle{plain}

\begin{center}
    \hypertarget{sec:toc}{}%
    {\Large \textbf{Contents}}\\[0.5cm]
\end{center}

\setlength{\parindent}{0pt}
\setlength{\parskip}{3pt}
\justifying

\noindent \hyperlink{sec:declaration}{Declaration \dotfill \textit{i}}\\
\hyperlink{sec:certificate}{Certificate \dotfill \textit{ii}}\\
\hyperlink{sec:abstract}{Abstract \dotfill \textit{iii}}\\
\hyperlink{sec:ack}{Acknowledgements \dotfill \textit{iv}}\\
\hyperlink{sec:toc}{Contents \dotfill \textit{v}}\\[0.5cm]

\hyperlink{sec:ch1}{\textbf{1 \hspace{0.3cm} Introduction} \dotfill \textit{2}}\\
\hspace*{0.8cm}\hyperlink{sec:s11}{1.1 \hspace{0.3cm} Background \dotfill \textit{2}}\\
\hspace*{0.8cm}\hyperlink{sec:s12}{1.2 \hspace{0.3cm} Motivation \dotfill \textit{2}}\\
\hspace*{0.8cm}\hyperlink{sec:s13}{1.3 \hspace{0.3cm} Objectives \dotfill \textit{3}}\\
\hspace*{0.8cm}\hyperlink{sec:s14}{1.4 \hspace{0.3cm} Scope and Contributions \dotfill \textit{4}}\\[0.3cm]

\hyperlink{sec:ch2}{\textbf{2 \hspace{0.3cm} Literature Review} \dotfill \textit{5}}\\
\hspace*{0.8cm}\hyperlink{sec:s21}{2.1 \hspace{0.3cm} The NTLSRU-Net Framework \dotfill \textit{5}}\\
\hspace*{0.8cm}\hyperlink{sec:s22}{2.2 \hspace{0.3cm} Background Context \dotfill \textit{5}}\\[0.3cm]

\hyperlink{sec:ch3}{\textbf{3 \hspace{0.3cm} Data Acquisition and Preprocessing} \dotfill \textit{7}}\\
\hspace*{0.8cm}\hyperlink{sec:s31}{3.1 \hspace{0.3cm} Data Source \dotfill \textit{7}}\\
\hspace*{0.8cm}\hyperlink{sec:s32}{3.2 \hspace{0.3cm} Study Areas \dotfill \textit{7}}\\
\hspace*{0.8cm}\hyperlink{sec:s33}{3.3 \hspace{0.3cm} Preprocessing Pipeline \dotfill \textit{8}}\\
\hspace*{0.8cm}\hyperlink{sec:s34}{3.4 \hspace{0.3cm} Train / Test / Validation Split \dotfill \textit{10}}\\[0.3cm]

\hyperlink{sec:ch4}{\textbf{4 \hspace{0.3cm} Methodology} \dotfill \textit{11}}\\
\hspace*{0.8cm}\hyperlink{sec:s41}{4.1 \hspace{0.3cm} SARIMA -- Statistical Baseline \dotfill \textit{11}}\\
\hspace*{0.8cm}\hyperlink{sec:s42}{4.2 \hspace{0.3cm} LSTM -- Deep Learning Scalar Model \dotfill \textit{15}}\\
\hspace*{0.8cm}\hyperlink{sec:s43}{4.3 \hspace{0.3cm} ConvLSTM -- Spatial Deep Learning Model \dotfill \textit{17}}\\[0.3cm]

\hyperlink{sec:ch5}{\textbf{5 \hspace{0.3cm} System Architecture and Dashboard} \dotfill \textit{22}}\\
\hspace*{0.8cm}\hyperlink{sec:s51}{5.1 \hspace{0.3cm} Overall Pipeline Architecture \dotfill \textit{22}}\\
\hspace*{0.8cm}\hyperlink{sec:s52}{5.2 \hspace{0.3cm} Codebase Structure \dotfill \textit{23}}\\
\hspace*{0.8cm}\hyperlink{sec:s53}{5.3 \hspace{0.3cm} Dashboard -- Lumina \dotfill \textit{24}}\\
\hspace*{0.8cm}\hyperlink{sec:s54}{5.4 \hspace{0.3cm} Deployment \dotfill \textit{25}}\\[0.3cm]

\hyperlink{sec:ch6}{\textbf{6 \hspace{0.3cm} Results and Discussion} \dotfill \textit{27}}\\
\hspace*{0.8cm}\hyperlink{sec:s61}{6.1 \hspace{0.3cm} Model Comparison -- Kharagpur \dotfill \textit{27}}\\
\hspace*{0.8cm}\hyperlink{sec:s62}{6.2 \hspace{0.3cm} Model Comparison -- Kolkata \dotfill \textit{28}}\\
\hspace*{0.8cm}\hyperlink{sec:s63}{6.3 \hspace{0.3cm} Discussion \dotfill \textit{28}}\\[0.3cm]

\hyperlink{sec:ch7}{\textbf{7 \hspace{0.3cm} Conclusion and Future Work} \dotfill \textit{30}}\\
\hspace*{0.8cm}\hyperlink{sec:s71}{7.1 \hspace{0.3cm} Conclusion \dotfill \textit{30}}\\
\hspace*{0.8cm}\hyperlink{sec:s72}{7.2 \hspace{0.3cm} Future Work -- Near Term \dotfill \textit{30}}\\
\hspace*{0.8cm}\hyperlink{sec:s73}{7.3 \hspace{0.3cm} Future Work -- Research Paper Proposal \dotfill \textit{31}}\\[0.3cm]

\hyperlink{sec:bib}{Bibliography \dotfill \textit{33}}\\

\pagenumbering{arabic}
\setcounter{page}{1}

% =====================================================================
%  CHAPTER 1 — INTRODUCTION
% =====================================================================
\newpage
\thispagestyle{plain}
\begin{center}
    \hypertarget{sec:ch1}{}%
    {\Large \textbf{1. Introduction}}\\[0.3cm]
    \noindent\rule{\linewidth}{0.6pt}\\[0.8cm]
\end{center}
\setlength{\parindent}{1cm}
\setlength{\parskip}{6pt}
\justifying

Night-time light (NTL) imagery captured from Earth-observing satellites has emerged as one of the most directly interpretable proxies for human activity at a planetary scale. The glow of cities, industrial zones, and transportation corridors observed from space reflects electricity consumption, urban density, and economic output in a way that is independent of administrative boundaries and is freely available over multi-decadal timescales. The ability to not merely observe historical NTL patterns but to \emph{forecast} future NTL values has broad implications for urban planning, energy policy, disaster response, and economic monitoring.

This report presents \textbf{Lumina} --- a multi-model nighttime light forecasting platform developed as BTP-2 at IIT Kharagpur. Lumina implements three forecasting models of increasing spatial complexity, validates them against real 2026 satellite data, and exposes all results through a publicly accessible interactive dashboard.

% ----- 1.1 ------------------------------------------------------------------
\vspace{0.5cm}
\begin{center}
    \hypertarget{sec:s11}{}%
    {\large \textbf{1.1 Background}}\\[0.3cm]
    \noindent\rule{\linewidth}{0.4pt}\\[0.5cm]
\end{center}

The primary satellite instrument used in this work is the \textbf{Visible Infrared Imaging Radiometer Suite} (VIIRS), launched aboard the Suomi-NPP satellite in 2011 and subsequently NOAA-20. VIIRS produces calibrated, cloud-free monthly nighttime radiance composites at 500 metres per pixel. The band used throughout this project, \texttt{avg\_rad}, is a physically meaningful measure of emitted light intensity in units of nW/cm$^2$/sr (nanowatts per square centimetre per steradian), corrected for stray light, lightning, and cloud cover.

Prior to VIIRS, the Defence Meteorological Satellite Programme's Operational Linescan System (DMSP-OLS) provided nighttime imagery from 1992 onwards, but DMSP data is uncalibrated and prone to saturation over bright urban areas. Chen et al.\ \cite{chen2024} proposed \textbf{NTLSRU-Net}, a U-Net-based deep learning framework that reconstructs VIIRS-like radiance from DMSP inputs, providing a continuous calibrated record from 1992 to the present. BTP-1 implemented and validated this reconstruction pipeline for an Indian region (Buxar, Bihar, 2012). BTP-2 builds directly on that foundation, taking the calibrated VIIRS time series as its starting point and asking: \emph{given a decade of monthly NTL data, can we reliably forecast future values?}

% ----- 1.2 ------------------------------------------------------------------
\vspace{0.5cm}
\begin{center}
    \hypertarget{sec:s12}{}%
    {\large \textbf{1.2 Motivation}}\\[0.3cm]
    \noindent\rule{\linewidth}{0.4pt}\\[0.5cm]
\end{center}

Several converging factors motivate the development of an NTL forecasting platform:

\begin{itemize}[leftmargin=1.5cm]
    \item \textbf{Urban growth planning:} Municipal authorities can use NTL forecasts to anticipate the spatial expansion of electrified zones and plan infrastructure investment accordingly.
    \item \textbf{Energy demand estimation:} NTL radiance correlates with electricity consumption. Short-range NTL forecasts can serve as auxiliary inputs to energy grid management systems.
    \item \textbf{Disaster recovery monitoring:} After natural disasters or conflicts, the trajectory of NTL recovery reveals how quickly affected areas return to pre-event activity levels. Forecasts provide a counterfactual baseline for comparison.
    \item \textbf{Economic proxies:} District-level NTL growth tracks GDP growth in data-scarce regions where official statistics are infrequent or unreliable.
    \item \textbf{Multi-modal research foundation:} Forecasted NTL series, when combined with population projections and GDP estimates, can serve as inputs to richer multi-modal models --- a direction this project proposes for future research.
\end{itemize}

At the time of this project, no publicly accessible platform existed that demonstrated multi-model NTL forecasting for Indian cities, with live validation against satellite-confirmed 2026 observations.

% ----- 1.3 ------------------------------------------------------------------
\vspace{0.5cm}
\begin{center}
    \hypertarget{sec:s13}{}%
    {\large \textbf{1.3 Objectives}}\\[0.3cm]
    \noindent\rule{\linewidth}{0.4pt}\\[0.5cm]
\end{center}

The specific objectives of this project are:

\begin{enumerate}[leftmargin=1.5cm]
    \item Download and preprocess monthly VIIRS VCMSLCFG nighttime light data for two Indian cities --- Kharagpur and Kolkata --- covering the period January 2014 to March 2026.
    \item Implement and evaluate three forecasting models spanning statistical and deep learning paradigms:
        \begin{itemize}
            \item SARIMA: a seasonal ARIMA statistical model as the baseline.
            \item LSTM: a stacked recurrent neural network with MC Dropout uncertainty estimation.
            \item ConvLSTM: a spatial encoder-decoder for pixel-level nighttime light forecasting.
        \end{itemize}
    \item Validate all forecasts against real January--March 2026 VIIRS observations (live validation), providing a rigorous evaluation beyond held-out test sets.
    \item Build an interactive, multi-page web dashboard (\textbf{Lumina}) that exposes all results, raw satellite maps, change detection, and model comparisons to end users.
    \item Identify the best-performing model and propose a concrete roadmap for future multi-modal NTL research.
\end{enumerate}

% ----- 1.4 ------------------------------------------------------------------
\vspace{0.5cm}
\begin{center}
    \hypertarget{sec:s14}{}%
    {\large \textbf{1.4 Scope and Contributions}}\\[0.3cm]
    \noindent\rule{\linewidth}{0.4pt}\\[0.5cm]
\end{center}

The scope of this project is limited to two West Bengal cities (Kharagpur and Kolkata) and the VIIRS \texttt{avg\_rad} band. The primary contributions of this work are:

\begin{itemize}[leftmargin=1.5cm]
    \item An end-to-end, reproducible pipeline from raw GEE export to trained forecast models for multiple cities, controlled by a single city-registry configuration file.
    \item A three-way model comparison (SARIMA vs LSTM vs ConvLSTM) under identical data splits and evaluation conditions, including live 2026 holdout validation.
    \item The first pixel-level NTL forecasting implementation for these Indian cities using a ConvLSTM encoder-decoder, yielding spatially explicit predictions suitable for change detection.
    \item A publicly deployed, open-source dashboard (\url{https://lumina-iamabhinav-dev.streamlit.app/}) providing interactive access to all results.
    \item A forward-looking research proposal integrated into the platform architecture, enabling future extension to multi-modal (NTL + population + GDP) analytics.
\end{itemize}

% =====================================================================
%  CHAPTER 2 — LITERATURE REVIEW
% =====================================================================
\newpage
\thispagestyle{plain}
\begin{center}
    \hypertarget{sec:ch2}{}%
    {\Large \textbf{2. Literature Review}}\\[0.3cm]
    \noindent\rule{\linewidth}{0.6pt}\\[0.8cm]
\end{center}
\setlength{\parindent}{1cm}
\setlength{\parskip}{6pt}
\justifying

This chapter covers the one formally cited paper on which this project directly builds~\cite{chen2024}, followed by descriptive background on the broader methodological context.

% ----- 2.1 ------------------------------------------------------------------
\vspace{0.5cm}
\begin{center}
    \hypertarget{sec:s21}{}%
    {\large \textbf{2.1 The NTLSRU-Net Framework}}\\[0.3cm]
    \noindent\rule{\linewidth}{0.4pt}\\[0.5cm]
\end{center}

Chen et al.\ \cite{chen2024} presented a global annual simulated VIIRS nighttime light dataset covering 1992 to 2023, generated using a deep learning reconstruction framework called \textbf{NTLSRU-Net}. The central problem addressed by the paper is the calibration gap between two generations of NTL instruments: the legacy DMSP-OLS sensor (1992--2013), which produces saturated, uncalibrated radiance values, and VIIRS (2012--present), which is fully calibrated and much higher in dynamic range and spatial resolution. The absence of a continuous, calibrated NTL time series spanning both eras is a major barrier for long-term socioeconomic studies.

The proposed architecture is a U-Net with skip connections, taking two input channels --- a DMSP-OLS NTL layer and a Normalised Difference Vegetation Index (NDVI) layer --- and producing a single-channel VIIRS-like radiance output at 500 metre resolution. The model was trained on paired DMSP-VIIRS observations from 2012--2013, the overlap years in which both sensors were operating simultaneously. Pre-trained weights are made publicly available. The paper validated the reconstruction using $R^2$ and RMSE across diverse land cover types and multiple world regions, demonstrating high-fidelity reconstruction of historical NTL patterns.

\textbf{Relevance to this project.} BTP-1 reproduced and applied the NTLSRU-Net pipeline to reconstruct a calibrated NTL time series for an Indian region (Buxar, Bihar). BTP-2 takes the \emph{output} of this pipeline as its starting point: a clean, calibrated VIIRS-like monthly time series. The existence of a reliable long-term NTL record --- made possible by the Chen et al.\ framework --- is the direct prerequisite for building the multi-year forecasting models implemented in this project.

% ----- 2.2 ------------------------------------------------------------------
\vspace{0.5cm}
\begin{center}
    \hypertarget{sec:s22}{}%
    {\large \textbf{2.2 Background Context}}\\[0.3cm]
    \noindent\rule{\linewidth}{0.4pt}\\[0.5cm]
\end{center}

\textbf{NTL as a socioeconomic proxy.} Nighttime light data has long been used as a proxy indicator for a range of socioeconomic phenomena including urbanisation, electrification progress, economic activity, and population density. The annual and seasonal patterns visible in NTL time series reflect agricultural cycles, festive seasons, monsoon-driven weather, and the long-term trajectory of urban growth. These characteristics make NTL a rich target for time-series forecasting.

\textbf{SARIMA.} The Seasonal Autoregressive Integrated Moving Average (SARIMA) model is a classical Box-Jenkins approach to univariate time series forecasting. It explicitly models autocorrelation structure, seasonal periodicity, and non-stationarity via differencing. SARIMA has been a standard baseline in meteorological and environmental time-series literature for decades. Its computational cost and data requirements are low, making it particularly competitive when the time series is short (100--200 observations) relative to the complexity of a deep learning model.

\textbf{LSTM.} Long Short-Term Memory (LSTM) networks are a recurrent neural network variant that address the vanishing gradient problem through input, forget, and output gating mechanisms. They have been widely applied to forecasting tasks including weather prediction, traffic flow, and energy demand estimation. The sliding-window formulation used in this project --- where the network is trained to map a window of $w$ past observations to the next value --- is a standard sequence-to-scalar regression approach.

\textbf{ConvLSTM.} ConvLSTM2D, introduced for precipitation nowcasting, replaces the linear transformations in LSTM gates with spatial convolutions, enabling the model to operate on sequences of spatial grids rather than scalar values. This makes it naturally suited to satellite image sequence forecasting where spatial correlation between neighbouring pixels carries predictive information. The encoder-decoder architecture used in this project encodes a window of historical frames into a latent spatial-temporal representation, then decodes it into a single predicted future frame.

\textbf{Monte Carlo Dropout.} Monte Carlo (MC) Dropout is a Bayesian approximation technique that enables uncertainty quantification in neural networks without any architectural changes beyond regular dropout. By performing multiple stochastic forward passes at inference time (with dropout active), the distribution of outputs approximates a posterior predictive distribution, with the mean as the point forecast and the standard deviation as the uncertainty estimate.

% =====================================================================
%  CHAPTER 3 — DATA ACQUISITION AND PREPROCESSING
% =====================================================================
\newpage
\thispagestyle{plain}
\begin{center}
    \hypertarget{sec:ch3}{}%
    {\Large \textbf{3. Data Acquisition and Preprocessing}}\\[0.3cm]
    \noindent\rule{\linewidth}{0.6pt}\\[0.8cm]
\end{center}
\setlength{\parindent}{1cm}
\setlength{\parskip}{6pt}
\justifying

This chapter describes the satellite data source, the two study areas selected for analysis, the end-to-end preprocessing pipeline, and the temporal split used for model training, testing, and live validation.

% ----- 3.1 ------------------------------------------------------------------
\vspace{0.5cm}
\begin{center}
    \hypertarget{sec:s31}{}%
    {\large \textbf{3.1 Data Source}}\\[0.3cm]
    \noindent\rule{\linewidth}{0.4pt}\\[0.5cm]
\end{center}

The primary data source for this project is the \textbf{VIIRS VCMSLCFG} (Visible Infrared Imaging Radiometer Suite -- Version 1 Cloud-Free Monthly Composites) product, made available via the \textbf{Google Earth Engine} (GEE) public data catalogue. VIIRS is a sensor aboard the Suomi-NPP and NOAA-20 satellites launched in 2011, providing calibrated, cloud-free monthly composites of nighttime radiance at a ground resolution of \textbf{500 metres per pixel}.

The band used throughout this project is \texttt{avg\_rad}, which records average nighttime radiance in units of \textbf{nW/cm$^2$/sr} (nanowatts per square centimetre per steradian). This is a physically calibrated, stable measure of emitted light intensity, making it suitable for time-series analysis across multiple years.

Data was downloaded for the period \textbf{January 2014 to March 2026} (147 monthly composite GeoTIFF images) for each of the two study cities. Download was automated using the GEE Python API via the script \texttt{src/download\_data.py}.

% ----- 3.2 ------------------------------------------------------------------
\vspace{0.5cm}
\begin{center}
    \hypertarget{sec:s32}{}%
    {\large \textbf{3.2 Study Areas}}\\[0.3cm]
    \noindent\rule{\linewidth}{0.4pt}\\[0.5cm]
\end{center}

Two cities in the Indian state of West Bengal were selected as study areas representing two distinct contexts: a small campus city and a major metropolitan centre.

\begin{table}[H]
\centering
\caption*{\textbf{Table 3.1:} Study area bounding boxes and grid properties.}
\begin{tabular}{lll}
\toprule
\textbf{Property} & \textbf{Kharagpur} & \textbf{Kolkata} \\
\midrule
Longitude range   & 87.25°E -- 87.45°E & 88.20°E -- 88.55°E \\
Latitude range    & 22.30°N -- 22.45°N & 22.40°N -- 22.75°N \\
Approximate grid  & ${\sim}35 \times 45$ pixels & ${\sim}70 \times 70$ pixels \\
Pixels per frame  & ${\sim}1{,}575$ & ${\sim}4{,}900$ \\
GEE scale         & 500 m/pixel & 500 m/pixel \\
Character         & Small city / campus & Major metro \\
State             & West Bengal, India & West Bengal, India \\
\bottomrule
\end{tabular}
\end{table}

\textbf{Kharagpur} is a small city dominated by the IIT Kharagpur campus and associated railway infrastructure. Its NTL signal is relatively low in absolute magnitude and exhibits a clear annual seasonality tied to the academic calendar and weather patterns. The compact spatial extent (35$\times$45 pixels) makes it computationally tractable for the ConvLSTM spatial model.

\textbf{Kolkata} is one of India's largest cities with a population exceeding 14 million in the metropolitan area. Its NTL signal is substantially higher in absolute radiance, exhibits more complex spatial heterogeneity, and presents a more challenging forecasting problem due to a noisier seasonal pattern.

% ----- 3.3 ------------------------------------------------------------------
\vspace{0.5cm}
\begin{center}
    \hypertarget{sec:s33}{}%
    {\large \textbf{3.3 Preprocessing Pipeline}}\\[0.3cm]
    \noindent\rule{\linewidth}{0.4pt}\\[0.5cm]
\end{center}

Each downloaded GeoTIFF contains one month of cloud-free VIIRS nighttime radiance for the study area bounding box. The preprocessing pipeline proceeds in the following steps, implemented in \texttt{src/preprocess.py} and \texttt{src/download\_data.py}:

\begin{enumerate}[leftmargin=1.5cm]
    \item \textbf{Nodata masking:} Pixels with \texttt{avg\_rad} $> 1{,}000$ nW/cm$^2$/sr are treated as sensor fill values or cloud artefacts and set to \texttt{NaN}. Values at or below the threshold are retained.

    \item \textbf{Mean brightness extraction:} For each monthly frame, the spatial mean of all valid (non-NaN) pixels is computed using \texttt{numpy.nanmean}. This produces a single scalar representing the mean nighttime radiance over the study area for that month.

    \item \textbf{Outlier cleaning:} The mean radiance time series is inspected for anomalous spikes (months where radiance is implausibly high relative to surrounding months). Flagged outliers are replaced with a locally smoothed value, producing the cleaned series \texttt{mean\_brightness\_clean.csv}.

    \item \textbf{Frame stacking for ConvLSTM:} The full spatial grid for each month is stored as a 3-D array $(H \times W \times 1)$. All monthly frames are stacked into a 4-D array and saved as \texttt{frames.npz} together with a \texttt{frame\_metadata.json} file containing dates, bounding box, and per-pixel statistics.
\end{enumerate}

\begin{figure}[H]
    \centering
    \includegraphics[width=0.85\textwidth]{kharagpur/sarima/eda_01_timeseries.png}
    \caption*{\textbf{Figure 3.1:} Raw mean NTL time series for Kharagpur (January 2014 -- March 2026). The annual seasonality and a weak upward trend are visible.}
\end{figure}

\begin{figure}[H]
    \centering
    \includegraphics[width=0.85\textwidth]{kolkata/sarima/eda_01_timeseries.png}
    \caption*{\textbf{Figure 3.2:} Raw mean NTL time series for Kolkata (January 2014 -- March 2026). The signal is substantially higher in absolute magnitude and shows a stronger long-term growth trend.}
\end{figure}

\begin{figure}[H]
    \centering
    \includegraphics[width=0.85\textwidth]{kharagpur/sarima/cleaning_report.png}
    \caption*{\textbf{Figure 3.3:} Data cleaning summary for Kharagpur showing original vs cleaned mean radiance and flagged outlier months.}
\end{figure}

% ----- 3.4 ------------------------------------------------------------------
\vspace{0.5cm}
\begin{center}
    \hypertarget{sec:s34}{}%
    {\large \textbf{3.4 Train / Test / Validation Split}}\\[0.3cm]
    \noindent\rule{\linewidth}{0.4pt}\\[0.5cm]
\end{center}

The 147-month time series (January 2014 -- March 2026) was partitioned into three non-overlapping segments:

\begin{table}[H]
\centering
\caption*{\textbf{Table 3.2:} Temporal data split used across all models.}
\begin{tabular}{llll}
\toprule
\textbf{Split} & \textbf{Period} & \textbf{Months} & \textbf{Purpose} \\
\midrule
Training  & Jan 2014 -- Dec 2023 & 120 & Model fitting \\
Test      & Jan 2024 -- Dec 2025 & 24  & In-sample evaluation (held-out) \\
Live validation & Jan 2026 -- Mar 2026 & 3  & Out-of-sample 2026 actual data \\
\bottomrule
\end{tabular}
\end{table}

All models are trained exclusively on the 120-month training set. Hyperparameter selection (for SARIMA order and LSTM architecture) uses only training data. The 24-month test set is used solely for final evaluation metrics. The three-month live validation set (January--March 2026) was downloaded in April 2026 and represents fully unseen future observations at the time the models were frozen.

\begin{figure}[H]
    \centering
    \includegraphics[width=0.85\textwidth]{kharagpur/sarima/train_test_split.png}
    \caption*{\textbf{Figure 3.4:} Train/test split visualisation for Kharagpur showing the 120-month training period and the 24-month held-out test period.}
\end{figure}

% =====================================================================
%  CHAPTER 4 — METHODOLOGY
% =====================================================================
\newpage
\thispagestyle{plain}
\begin{center}
    \hypertarget{sec:ch4}{}%
    {\Large \textbf{4. Methodology}}\\[0.3cm]
    \noindent\rule{\linewidth}{0.6pt}\\[0.8cm]
\end{center}
\setlength{\parindent}{1cm}
\setlength{\parskip}{6pt}
\justifying

Three forecasting models were implemented, each addressing a different modelling paradigm: a classical statistical model (SARIMA), a scalar deep-learning recurrent model (LSTM), and a spatial deep-learning model (ConvLSTM). This section describes the design, training procedure, and architecture of each.

% ----- 4.1 ------------------------------------------------------------------
\vspace{0.5cm}
\begin{center}
    \hypertarget{sec:s41}{}%
    {\large \textbf{4.1 SARIMA -- Statistical Baseline}}\\[0.3cm]
    \noindent\rule{\linewidth}{0.4pt}\\[0.5cm]
\end{center}

\textbf{Model description.} SARIMA (Seasonal AutoRegressive Integrated Moving Average) is a classical Box-Jenkins family model for univariate time series with trend and seasonality. It is specified by two sets of parameters: the non-seasonal order $(p, d, q)$ controlling the autoregressive, differencing, and moving average components, and the seasonal order $(P, D, Q)[s]$ with period $s = 12$ (months). The model equation for a doubly-differenced series $w_t$ is:

\begin{equation}
    \Phi_P(B^{12})\,\phi_p(B)\,w_t = \Theta_Q(B^{12})\,\theta_q(B)\,\varepsilon_t, \quad \varepsilon_t \sim \mathcal{N}(0, \sigma^2)
\end{equation}

where $B$ is the backshift operator, $\phi_p$ and $\Phi_P$ are the non-seasonal and seasonal AR polynomials, $\theta_q$ and $\Theta_Q$ are the corresponding MA polynomials, and $w_t = (1 - B)^d (1 - B^{12})^D x_t$.

\textbf{Order selection.} The order was selected by exhaustive grid search over all combinations of $p \in \{0,1,2\}$, $d \in \{0,1\}$, $q \in \{0,1,2\}$, $P \in \{0,1\}$, $D \in \{0,1\}$, $Q \in \{0,1\}$ with $s = 12$, minimising Akaike Information Criterion (AIC) on the training set. The search identified:

\begin{table}[H]
\centering
\caption*{\textbf{Table 4.1:} SARIMA best orders by city (selected by AIC on training data).}
\begin{tabular}{lll}
\toprule
\textbf{City} & \textbf{Best Order} & \textbf{Training AIC} \\
\midrule
Kharagpur & SARIMA$(0,1,1)(0,1,1)_{12}$ & 189.95 \\
Kolkata   & SARIMA$(0,1,2)(0,1,1)_{12}$ & 444.22 \\
\bottomrule
\end{tabular}
\end{table}

The selected orders indicate that one regular difference and one seasonal difference are sufficient to achieve stationarity in both cities, consistent with the upward trend and strong annual periodicity visible in the time series.

\textbf{Stationarity analysis.} Before fitting, ACF and PACF plots were inspected at three stages: raw series, after first difference, and after first seasonal difference. These plots guided confirmation of the final order selection.

\begin{figure}[H]
    \centering
    \includegraphics[width=0.92\textwidth]{kharagpur/sarima/eda_05_acf_pacf_raw.png}
    \caption*{\textbf{Figure 4.1:} ACF and PACF of the raw Kharagpur time series showing significant autocorrelation at lags 12, 24 confirming the seasonal period.}
\end{figure}

\begin{figure}[H]
    \centering
    \includegraphics[width=0.92\textwidth]{kharagpur/sarima/eda_06_acf_pacf_diff1.png}
    \caption*{\textbf{Figure 4.2:} ACF and PACF after first-order differencing of the Kharagpur series: trend is removed but seasonal structure remains.}
\end{figure}

\begin{figure}[H]
    \centering
    \includegraphics[width=0.92\textwidth]{kharagpur/sarima/eda_07_acf_pacf_diff1_s1.png}
    \caption*{\textbf{Figure 4.3:} ACF and PACF after first-order + seasonal differencing: the series is near-stationary, justifying the $(d=1, D=1)$ order.}
\end{figure}

\textbf{Supplementary EDA.} Seasonal boxplots and year-over-year overlays were produced to visualise further characteristics of the time series.

\begin{figure}[H]
    \centering
    \begin{minipage}[t]{0.48\textwidth}
        \centering
        \includegraphics[width=\textwidth]{kharagpur/sarima/eda_02_seasonal_boxplot.png}
        \caption*{\textbf{Figure 4.4:} Monthly seasonal boxplot (Kharagpur).}
    \end{minipage}
    \hfill
    \begin{minipage}[t]{0.48\textwidth}
        \centering
        \includegraphics[width=\textwidth]{kharagpur/sarima/eda_03_yoy_overlay.png}
        \caption*{\textbf{Figure 4.5:} Year-over-year NTL overlay (Kharagpur).}
    \end{minipage}
\end{figure}

\textbf{Training.} The model was fit using the \texttt{statsmodels} Python library on the 120-month training series (January 2014 -- December 2023).

\textbf{Forecasting.} The fitted model produces point forecasts with 95\% confidence intervals. For the test evaluation, the model forecasts 24 months ahead (January 2024 -- December 2025). For the 2026 live validation, the model was re-applied to produce January--March 2026 forecasts.

\textbf{Test-set results.}

\begin{figure}[H]
    \centering
    \includegraphics[width=0.90\textwidth]{kharagpur/sarima/eval_01_forecast_vs_actual.png}
    \caption*{\textbf{Figure 4.6:} SARIMA test-set forecast vs actual for Kharagpur (Jan 2024 -- Dec 2025) with 95\% confidence interval.}
\end{figure}

\begin{figure}[H]
    \centering
    \includegraphics[width=0.90\textwidth]{kharagpur/sarima/eval_05_diagnostics_panel.png}
    \caption*{\textbf{Figure 4.7:} SARIMA residual diagnostics for Kharagpur: histogram, Q-Q plot, ACF of residuals, and standardized residuals over time.}
\end{figure}

\begin{figure}[H]
    \centering
    \includegraphics[width=0.90\textwidth]{kolkata/sarima/eval_01_forecast_vs_actual.png}
    \caption*{\textbf{Figure 4.8:} SARIMA test-set forecast vs actual for Kolkata (Jan 2024 -- Dec 2025).}
\end{figure}

\begin{table}[H]
\centering
\caption*{\textbf{Table 4.2:} SARIMA test-set evaluation metrics (Jan 2024 -- Dec 2025).}
\begin{tabular}{lcccc}
\toprule
\textbf{City} & \textbf{MAE} & \textbf{RMSE} & \textbf{MAPE (\%)} & \textbf{MASE} \\
\midrule
Kharagpur & 0.838 & 1.295 & 13.75 & 1.260 \\
Kolkata   & 2.545 & 4.052 & 24.82 & 1.151 \\
\bottomrule
\end{tabular}
\end{table}

\textbf{2026 forecast.}

\begin{figure}[H]
    \centering
    \begin{minipage}[t]{0.48\textwidth}
        \centering
        \includegraphics[width=\textwidth]{kharagpur/sarima/forecast.png}
        \caption*{\textbf{Figure 4.9:} SARIMA 2026 forecast with 95\% CI (Kharagpur).}
    \end{minipage}
    \hfill
    \begin{minipage}[t]{0.48\textwidth}
        \centering
        \includegraphics[width=\textwidth]{kolkata/sarima/forecast.png}
        \caption*{\textbf{Figure 4.10:} SARIMA 2026 forecast with 95\% CI (Kolkata).}
    \end{minipage}
\end{figure}

% ----- 4.2 ------------------------------------------------------------------
\vspace{0.5cm}
\begin{center}
    \hypertarget{sec:s42}{}%
    {\large \textbf{4.2 LSTM -- Deep Learning Scalar Model}}\\[0.3cm]
    \noindent\rule{\linewidth}{0.4pt}\\[0.5cm]
\end{center}

\textbf{Model description.} A stacked Long Short-Term Memory (LSTM) network was used as the deep learning scalar baseline. LSTM is a recurrent neural network architecture designed to capture long-range temporal dependencies through gated memory cells. The network takes a sliding window of $w$ past observations as input and predicts the next value.

\textbf{Architecture.} The final architecture (after hyperparameter tuning) is:

\begin{lstlisting}[caption=LSTM architecture (Keras pseudo-code)]
Input(shape=(window_size, 1))
LSTM(units_1, return_sequences=True)
Dropout(dropout)
LSTM(units_2, return_sequences=False)
Dropout(dropout)
Dense(16, activation='relu')
Dense(1)

Loss: MAE  |  Optimizer: Adam  |  Metric: RMSE
\end{lstlisting}

\textbf{Hyperparameter tuning.} Keras Tuner Random Search was used with 30 trials to optimise over: window size $w \in \{6, 12, 24\}$, first LSTM layer units $\in \{32, 64, 128\}$, second LSTM layer units $\in \{16, 32\}$, dropout rate $\in \{0.1, 0.2, 0.3\}$, and learning rate $\in \{0.001, 0.0005, 0.0001\}$. The best configurations found were:

\begin{table}[H]
\centering
\caption*{\textbf{Table 4.3:} LSTM hyperparameters selected by Keras Tuner Random Search (30 trials).}
\begin{tabular}{lll}
\toprule
\textbf{Hyperparameter} & \textbf{Kharagpur} & \textbf{Kolkata} \\
\midrule
Window size     & 12 & 24 \\
Units (LSTM-1)  & 128 & 128 \\
Units (LSTM-2)  & 32  & 16  \\
Dropout         & 0.3 & 0.2 \\
Learning rate   & 0.0005 & 0.0005 \\
\bottomrule
\end{tabular}
\end{table}

\textbf{Monte Carlo Dropout (MC Dropout).} To obtain predictive uncertainty estimates, dropout is activated at inference time. Fifty forward passes are performed for each test timestamp, yielding a distribution of predictions whose mean serves as the point estimate and whose standard deviation provides the uncertainty band.

\textbf{Training.}

\begin{figure}[H]
    \centering
    \begin{minipage}[t]{0.48\textwidth}
        \centering
        \includegraphics[width=\textwidth]{kharagpur/lstm/loss_curve.png}
        \caption*{\textbf{Figure 4.11:} LSTM training and validation loss curve (Kharagpur).}
    \end{minipage}
    \hfill
    \begin{minipage}[t]{0.48\textwidth}
        \centering
        \includegraphics[width=\textwidth]{kolkata/lstm/loss_curve.png}
        \caption*{\textbf{Figure 4.12:} LSTM training and validation loss curve (Kolkata).}
    \end{minipage}
\end{figure}

\textbf{Test-set results.}

\begin{figure}[H]
    \centering
    \includegraphics[width=0.90\textwidth]{kharagpur/lstm/eval_01_forecast_vs_actual.png}
    \caption*{\textbf{Figure 4.13:} LSTM test-set forecast vs actual for Kharagpur with MC Dropout uncertainty band (Jan 2024 -- Dec 2025).}
\end{figure}

\begin{figure}[H]
    \centering
    \includegraphics[width=0.90\textwidth]{kolkata/lstm/eval_01_forecast_vs_actual.png}
    \caption*{\textbf{Figure 4.14:} LSTM test-set forecast vs actual for Kolkata with MC Dropout uncertainty band (Jan 2024 -- Dec 2025).}
\end{figure}

\begin{table}[H]
\centering
\caption*{\textbf{Table 4.4:} LSTM test-set evaluation metrics (Jan 2024 -- Dec 2025).}
\begin{tabular}{lcccc}
\toprule
\textbf{City} & \textbf{MAE} & \textbf{RMSE} & \textbf{MAPE (\%)} & \textbf{MASE} \\
\midrule
Kharagpur & 0.981 & 1.390 & 17.50 & 1.440 \\
Kolkata   & 3.357 & 5.147 & 30.50 & 1.512 \\
\bottomrule
\end{tabular}
\end{table}

\textbf{2026 forecast.}

\begin{figure}[H]
    \centering
    \begin{minipage}[t]{0.48\textwidth}
        \centering
        \includegraphics[width=\textwidth]{kharagpur/lstm/forecast.png}
        \caption*{\textbf{Figure 4.15:} LSTM 2026 forecast with uncertainty band (Kharagpur).}
    \end{minipage}
    \hfill
    \begin{minipage}[t]{0.48\textwidth}
        \centering
        \includegraphics[width=\textwidth]{kolkata/lstm/forecast.png}
        \caption*{\textbf{Figure 4.16:} LSTM 2026 forecast with uncertainty band (Kolkata).}
    \end{minipage}
\end{figure}

% ----- 4.3 ------------------------------------------------------------------
\vspace{0.5cm}
\begin{center}
    \hypertarget{sec:s43}{}%
    {\large \textbf{4.3 ConvLSTM -- Spatial Deep Learning Model}}\\[0.3cm]
    \noindent\rule{\linewidth}{0.4pt}\\[0.5cm]
\end{center}

\textbf{Model description.} ConvLSTM2D (Convolutional LSTM) replaces the matrix multiplications in a standard LSTM cell with 2-D convolutions. This allows the model to operate directly on spatial image sequences, capturing both temporal dynamics and spatial correlations simultaneously. While SARIMA and LSTM predict a scalar mean radiance per month, the ConvLSTM predicts an entire spatial grid, enabling pixel-level forecasting of nighttime light.

\textbf{Architecture.} An encoder-decoder design was used:

\begin{lstlisting}[caption=ConvLSTM encoder-decoder architecture (Keras pseudo-code)]
Input: (batch, T=12, H, W, 1)

-- Encoder --
ConvLSTM2D(16 filters, 3x3, padding=same, return_seq=True)  -> BatchNorm
ConvLSTM2D(32 filters, 3x3, padding=same, return_seq=True)  -> BatchNorm
ConvLSTM2D(32 filters, 3x3, padding=same, return_seq=False) -> BatchNorm
     [output shape: (batch, H, W, 32)]

-- Bridge: Reshape (batch, 1, H, W, 32) --

-- Decoder --
ConvLSTM2D(32 filters, 3x3, padding=same, return_seq=True)  -> BatchNorm
ConvLSTM2D(16 filters, 3x3, padding=same, return_seq=True)  -> BatchNorm

-- Output head --
TimeDistributed(Conv2D(1, 1x1, sigmoid)) -> Reshape (batch, H, W, 1)

Loss: MAE  |  Optimizer: Adam  |  Metric: pixel-RMSE
\end{lstlisting}

The \texttt{padding=same} setting throughout ensures the spatial dimensions are preserved exactly ($35 \times 45$ for Kharagpur, $70 \times 70$ for Kolkata). The sigmoid output head maps radiance values to the $[0, 1]$ normalised range (after per-pixel min-max normalisation at the frame-preparation stage).

\textbf{Training.}

\begin{figure}[H]
    \centering
    \begin{minipage}[t]{0.48\textwidth}
        \centering
        \includegraphics[width=\textwidth]{kharagpur/convlstm/loss_curve.png}
        \caption*{\textbf{Figure 4.17:} ConvLSTM training and validation loss (Kharagpur, 73 epochs).}
    \end{minipage}
    \hfill
    \begin{minipage}[t]{0.48\textwidth}
        \centering
        \includegraphics[width=\textwidth]{kolkata/convlstm/loss_curve.png}
        \caption*{\textbf{Figure 4.18:} ConvLSTM training and validation loss (Kolkata, 117 epochs).}
    \end{minipage}
\end{figure}

\textbf{Spatial evaluation metrics.} Because the ConvLSTM output is a 2-D image, different metrics are used:

\begin{itemize}[leftmargin=1.5cm]
    \item \textbf{Pixel-wise MAE / RMSE:} mean and root-mean-squared error computed over all spatial pixels and all test frames.
    \item \textbf{SSIM:} Structural Similarity Index Measure, which evaluates perceived image quality by comparing luminance, contrast, and structure in a patch-based manner.
    \item \textbf{PSNR:} Peak Signal-to-Noise Ratio (dB), quantifying the fidelity of the reconstructed frame relative to the actual.
    \item \textbf{Mean radiance MAE:} the MAE of the spatial mean of the predicted frame vs actual frame, enabling direct comparison with the scalar SARIMA and LSTM results.
\end{itemize}

\begin{table}[H]
\centering
\caption*{\textbf{Table 4.5:} ConvLSTM test-set evaluation metrics (Jan 2024 -- Dec 2025).}
\begin{tabular}{lcc}
\toprule
\textbf{Metric} & \textbf{Kharagpur} & \textbf{Kolkata} \\
\midrule
Pixel MAE       & 5.141 & 18.187 \\
Pixel RMSE      & 7.661 & 46.896 \\
Mean SSIM       & 0.477 & 0.536 \\
Mean PSNR (dB)  & 24.78 & 16.27 \\
Mean radiance MAE & 0.861 & 10.785 \\
\bottomrule
\end{tabular}
\end{table}

\textbf{Spatial results.}

\begin{figure}[H]
    \centering
    \includegraphics[width=0.92\textwidth]{kharagpur/convlstm/pred_vs_actual_grid.png}
    \caption*{\textbf{Figure 4.19:} ConvLSTM predicted vs actual spatial NTL frames for Kharagpur: sample months from the test period. Each pair shows actual (left) and predicted (right).}
\end{figure}

\begin{figure}[H]
    \centering
    \includegraphics[width=0.92\textwidth]{kharagpur/convlstm/error_maps.png}
    \caption*{\textbf{Figure 4.20:} Per-pixel mean absolute error map for Kharagpur. Brighter regions indicate consistent prediction errors.}
\end{figure}

\begin{figure}[H]
    \centering
    \includegraphics[width=0.90\textwidth]{kharagpur/convlstm/ssim_timeseries.png}
    \caption*{\textbf{Figure 4.21:} Per-month SSIM values for Kharagpur test period (Jan 2024 -- Dec 2025). Dips in SSIM correspond to rainy-season months with higher radiance variability.}
\end{figure}

\begin{figure}[H]
    \centering
    \includegraphics[width=0.92\textwidth]{kolkata/convlstm/pred_vs_actual_grid.png}
    \caption*{\textbf{Figure 4.22:} ConvLSTM predicted vs actual spatial NTL frames for Kolkata.}
\end{figure}

\begin{figure}[H]
    \centering
    \includegraphics[width=0.90\textwidth]{kolkata/convlstm/ssim_timeseries.png}
    \caption*{\textbf{Figure 4.23:} Per-month SSIM values for Kolkata test period.}
\end{figure}

\textbf{2026 spatial forecast.}

\begin{figure}[H]
    \centering
    \begin{minipage}[t]{0.48\textwidth}
        \centering
        \includegraphics[width=\textwidth]{kharagpur/convlstm/forecast_panel.png}
        \caption*{\textbf{Figure 4.24:} ConvLSTM 2026 spatial forecast panel (Kharagpur).}
    \end{minipage}
    \hfill
    \begin{minipage}[t]{0.48\textwidth}
        \centering
        \includegraphics[width=\textwidth]{kolkata/convlstm/forecast_panel.png}
        \caption*{\textbf{Figure 4.25:} ConvLSTM 2026 spatial forecast panel (Kolkata).}
    \end{minipage}
\end{figure}

% =====================================================================
%  CHAPTER 5 — SYSTEM ARCHITECTURE AND DASHBOARD
% =====================================================================
\newpage
\thispagestyle{plain}
\begin{center}
    \hypertarget{sec:ch5}{}%
    {\Large \textbf{5. System Architecture and Dashboard}}\\[0.3cm]
    \noindent\rule{\linewidth}{0.6pt}\\[0.8cm]
\end{center}
\setlength{\parindent}{1cm}
\setlength{\parskip}{6pt}
\justifying

This chapter describes the end-to-end software architecture of the Lumina project, the organisation of the codebase, the interactive dashboard built on top of the forecasting pipelines, and the deployment strategy.

% ----- 5.1 ------------------------------------------------------------------
\vspace{0.5cm}
\begin{center}
    \hypertarget{sec:s51}{}%
    {\large \textbf{5.1 Overall Pipeline Architecture}}\\[0.3cm]
    \noindent\rule{\linewidth}{0.4pt}\\[0.5cm]
\end{center}

The full system is organised as six sequential stages, each implemented as an independent module with clearly defined inputs and outputs:

\begin{verbatim}
  [1] DATA ACQUISITION
       src/download_data.py
       GEE Python API -> GeoTIFF per month -> data/{city}/tiffs/
               |
               v
  [2] PREPROCESSING
       src/preprocess.py
       GeoTIFF -> mask -> nanmean -> mean_brightness_clean.csv
       GeoTIFF -> stack -> frames.npz + frame_metadata.json
               |
               v
  [3] SARIMA PIPELINE
       models/sarima/clean_brightness.py  (cleaning)
       models/sarima/eda.py               (EDA + ACF/PACF plots)
       models/sarima/find_order.py        (grid search AIC)
       models/sarima/train.py             (fit model)
       models/sarima/evaluate.py          (test metrics + plots)
       models/sarima/forecast.py          (2026 forecast)
               |
               v
  [4] LSTM PIPELINE
       models/lstm/prepare_sequences.py   (sliding window)
       models/lstm/build_model.py         (architecture)
       models/lstm/tune.py                (Keras Tuner search)
       models/lstm/train.py               (final training)
       models/lstm/evaluate.py            (test metrics + plots)
       models/lstm/forecast.py            (2026 MC Dropout forecast)
               |
               v
  [5] ConvLSTM PIPELINE
       models/convlstm/prepare_frames.py  (frame normalisation)
       models/convlstm/build_convlstm.py  (architecture)
       models/convlstm/train_convlstm.py  (training)
       models/convlstm/evaluate_convlstm.py (spatial metrics)
       models/convlstm/forecast_convlstm.py (2026 spatial forecast)
               |
               v
  [6] DASHBOARD
       dashboard/app.py  +  dashboard/pages/0X_*.py
       Streamlit multi-page app -> lumina-iamabhinav-dev.streamlit.app
\end{verbatim}

Each stage reads its inputs from the previous stage's outputs directory and writes results to \texttt{outputs/{city}/{model}/}. All pipeline scripts are parameterised by city name (via \texttt{src/cities.py}), making it straightforward to add new cities without modifying the pipeline logic.

% ----- 5.2 ------------------------------------------------------------------
\vspace{0.5cm}
\begin{center}
    \hypertarget{sec:s52}{}%
    {\large \textbf{5.2 Codebase Structure}}\\[0.3cm]
    \noindent\rule{\linewidth}{0.4pt}\\[0.5cm]
\end{center}

The codebase follows a flat, module-based layout:

\begin{verbatim}
BTP/
|-- src/
|   |-- cities.py          City registry and path helpers
|   |-- download_data.py   GEE download automation
|   |-- preprocess.py      Raster loading, masking, brightness extraction
|   +-- utils.py           Shared utilities
|-- models/
|   |-- sarima/            SARIMA pipeline scripts
|   |-- lstm/              LSTM pipeline scripts
|   +-- convlstm/          ConvLSTM pipeline scripts
|       +-- {city}/        City-specific frame data
|-- outputs/
|   |-- sarima/            Kharagpur SARIMA results
|   |-- lstm/              Kharagpur LSTM results
|   |-- convlstm/          Kharagpur ConvLSTM results
|   +-- kolkata/           Kolkata results (sarima/, lstm/, convlstm/)
|-- data/
|   |-- tiffs/             Kharagpur GeoTIFFs
|   +-- kolkata/tiffs/     Kolkata GeoTIFFs
|-- dashboard/
|   |-- app.py             Streamlit app entry point
|   +-- pages/
|       |-- 01_explorer.py
|       |-- 02_change.py
|       |-- 03_charts.py
|       |-- 04_forecast.py
|       |-- 05_lstm.py
|       |-- 06_convlstm.py
|       +-- 07_compare.py
|-- requirements.txt       Python dependencies
+-- runtime.txt            Python version for Streamlit Cloud
\end{verbatim}

The city registry in \texttt{src/cities.py} stores each city's bounding box, GEE scale, display name, and chart colour. All pipeline scripts import this module to resolve city-specific paths, ensuring no hardcoded city assumptions appear anywhere in the pipeline code.

% ----- 5.3 ------------------------------------------------------------------
\vspace{0.5cm}
\begin{center}
    \hypertarget{sec:s53}{}%
    {\large \textbf{5.3 Dashboard -- Lumina}}\\[0.3cm]
    \noindent\rule{\linewidth}{0.4pt}\\[0.5cm]
\end{center}

\textbf{Lumina} is a multi-page Streamlit application that exposes all forecasting results through an interactive web interface. It is designed to be usable by non-technical stakeholders without any programming knowledge. The application consists of eight pages:

\begin{table}[H]
\centering
\caption*{\textbf{Table 5.1:} Dashboard pages and their purposes.}
\begin{tabular}{p{3.8cm}p{9.2cm}}
\toprule
\textbf{Page} & \textbf{Description} \\
\midrule
NTL Explorer & Interactive month browser with a time slider; renders false-colour VIIRS raster maps for any selected month and city. Supports animated playback over the full time series. \\
\addlinespace
Change Detection & Select two months and render a signed difference map showing areas of NTL increase (warm colours) and decrease (cool colours). \\
\addlinespace
Charts \& Trends & Time series plots, rolling mean, seasonal decomposition (trend, seasonality, residuals) for the selected city. \\
\addlinespace
SARIMA Forecast & Displays the fitted SARIMA model, test-period forecast vs actual, residual diagnostics, and the 2026 projection with confidence intervals. \\
\addlinespace
LSTM Forecast & MC Dropout forecast with uncertainty bands, training loss curve, and side-by-side comparison with SARIMA. \\
\addlinespace
ConvLSTM Spatial & Pixel-level spatial forecast maps, error heatmaps, SSIM time series, and GeoTIFF download for any forecast month. \\
\addlinespace
City Comparison & Side-by-side time series, seasonal patterns, and model performance metrics for Kharagpur vs Kolkata. \\
\addlinespace
2026 Validation & Live validation page: actual Jan--Mar 2026 data overlaid on all three model forecasts with per-month error percentages. \\
\bottomrule
\end{tabular}
\end{table}

All heavy computation (raster loading, model inference) is cached using Streamlit's \texttt{@st.cache\_data} and \texttt{@st.cache\_resource} decorators, ensuring fast page loads after the initial warm-up. Figures are rendered with Matplotlib and displayed inline via \texttt{st.pyplot()}.

% ----- 5.4 ------------------------------------------------------------------
\vspace{0.5cm}
\begin{center}
    \hypertarget{sec:s54}{}%
    {\large \textbf{5.4 Deployment}}\\[0.3cm]
    \noindent\rule{\linewidth}{0.4pt}\\[0.5cm]
\end{center}

The dashboard is deployed on \textbf{Streamlit Community Cloud} and is publicly accessible at:

\begin{center}
    \url{https://lumina-iamabhinav-dev.streamlit.app}
\end{center}

Deployment is managed via a GitHub repository (\texttt{iamabhinav-dev/Lumina}) connected to Streamlit Cloud. Every push to the \texttt{main} branch triggers an automatic redeploy. The Python version is pinned via \texttt{runtime.txt} (\texttt{python-3.11.x}) and all dependencies are specified in \texttt{requirements.txt}. Pre-computed outputs (trained model weights, forecast CSVs, evaluation metrics, plot PNGs) are committed to the repository so the deployed app does not need to retrain models at startup.

% =====================================================================
%  CHAPTER 6 — RESULTS AND DISCUSSION
% =====================================================================
\newpage
\thispagestyle{plain}
\begin{center}
    \hypertarget{sec:ch6}{}%
    {\Large \textbf{6. Results and Discussion}}\\[0.3cm]
    \noindent\rule{\linewidth}{0.6pt}\\[0.8cm]
\end{center}
\setlength{\parindent}{1cm}
\setlength{\parskip}{6pt}
\justifying

This chapter consolidates and compares the performance of all three models across both cities. Results are presented in two tiers: a standard held-out test evaluation (January 2024 -- December 2025, 24 months) and a live validation against actual January--March 2026 VIIRS composites downloaded in April 2026.

% ----- 6.1 ------------------------------------------------------------------
\vspace{0.5cm}
\begin{center}
    \hypertarget{sec:s61}{}%
    {\large \textbf{6.1 Model Comparison -- Kharagpur}}\\[0.3cm]
    \noindent\rule{\linewidth}{0.4pt}\\[0.5cm]
\end{center}

\textbf{Test-set performance (Jan 2024 -- Dec 2025).} All three models were evaluated on the 24-month held-out test set. For SARIMA and LSTM, scalar metrics (MAE, RMSE, MAPE, MASE) are computed on the mean brightness series. For ConvLSTM, the mean radiance MAE is included alongside spatial metrics to enable cross-model comparison.

\begin{table}[H]
\centering
\caption*{\textbf{Table 6.1:} Kharagpur test-set evaluation metrics (Jan 2024 -- Dec 2025).}
\begin{tabular}{lcccccc}
\toprule
\textbf{Model} & \textbf{MAE} & \textbf{RMSE} & \textbf{MAPE (\%)} & \textbf{MASE} & \textbf{SSIM} & \textbf{PSNR (dB)} \\
\midrule
SARIMA & 0.838 & 1.295 & 13.75 & 1.260 & --- & --- \\
LSTM   & 0.981 & 1.390 & 17.50 & 1.440 & --- & --- \\
ConvLSTM (mean rad) & 0.861 & 1.382 & --- & --- & 0.477 & 24.78 \\
\bottomrule
\end{tabular}
\end{table}

SARIMA achieves the lowest MAE (0.838) and MAPE (13.75\%) on the test set. LSTM is competitive but slightly weaker at MAE 0.981. ConvLSTM's mean-radiance MAE of 0.861 is close to SARIMA in scalar terms, confirming that its spatial prediction does not sacrifice average accuracy.

\textbf{Live validation against 2026 actual data.} The following table shows per-month actual VIIRS readings vs each model's predictions for January--March 2026.

\begin{table}[H]
\centering
\caption*{\textbf{Table 6.2:} Kharagpur live validation -- January--March 2026 (units: nW/cm$^2$/sr).}
\begin{tabular}{lccccccc}
\toprule
\textbf{Month} & \textbf{Actual} & \textbf{SARIMA} & \textbf{SARIMA err\%} & \textbf{LSTM} & \textbf{LSTM err\%} & \textbf{ConvLSTM} & \textbf{CLstm err\%} \\
\midrule
Jan 2026 & 9.154 & 8.182 & 10.6 & 7.617 & 16.8 & 7.743 & 15.4 \\
Feb 2026 & 9.462 & 8.486 & 10.3 & 7.390 & 21.9 & 7.738 & 18.2 \\
Mar 2026 & 8.869 & 8.058 & 9.1  & 7.648 & 13.8 & 7.726 & 12.9 \\
\midrule
\textbf{Mean} & --- & --- & \textbf{10.0} & --- & \textbf{17.5} & --- & \textbf{15.5} \\
\bottomrule
\end{tabular}
\end{table}

\textbf{Best model: SARIMA} with a mean live error of \textbf{10.0\%}. LSTM achieves 17.5\% and ConvLSTM achieves 15.5\%. All three models directionally follow the actual values (declining from February to March), confirming that the seasonal pattern is captured.

% ----- 6.2 ------------------------------------------------------------------
\vspace{0.5cm}
\begin{center}
    \hypertarget{sec:s62}{}%
    {\large \textbf{6.2 Model Comparison -- Kolkata}}\\[0.3cm]
    \noindent\rule{\linewidth}{0.4pt}\\[0.5cm]
\end{center}

\textbf{Test-set performance (Jan 2024 -- Dec 2025).}

\begin{table}[H]
\centering
\caption*{\textbf{Table 6.3:} Kolkata test-set evaluation metrics (Jan 2024 -- Dec 2025).}
\begin{tabular}{lcccccc}
\toprule
\textbf{Model} & \textbf{MAE} & \textbf{RMSE} & \textbf{MAPE (\%)} & \textbf{MASE} & \textbf{SSIM} & \textbf{PSNR (dB)} \\
\midrule
SARIMA & 2.545 & 4.052 & 24.82 & 1.151 & --- & --- \\
LSTM   & 3.357 & 5.147 & 30.50 & 1.512 & --- & --- \\
ConvLSTM (mean rad) & 10.785 & 12.546 & --- & --- & 0.536 & 16.27 \\
\bottomrule
\end{tabular}
\end{table}

For Kolkata, SARIMA again leads with MAE 2.545 and MAPE 24.82\%. The ConvLSTM's mean-radiance MAE is substantially higher (10.785) because the Kolkata spatial grid ($79 \times 79$ pixels) is more heterogeneous and the absolute radiance values are significantly larger.

\textbf{Live validation against 2026 actual data.}

\begin{table}[H]
\centering
\caption*{\textbf{Table 6.4:} Kolkata live validation -- January--March 2026 (units: nW/cm$^2$/sr).}
\begin{tabular}{lccccccc}
\toprule
\textbf{Month} & \textbf{Actual} & \textbf{SARIMA} & \textbf{SARIMA err\%} & \textbf{LSTM} & \textbf{LSTM err\%} & \textbf{ConvLSTM} & \textbf{CLstm err\%} \\
\midrule
Jan 2026 & 20.823 & 18.132 & 12.9 & 16.467 & 20.9 & 29.401 & 43.0 \\
Feb 2026 & 21.537 & 18.345 & 14.8 & 16.628 & 22.8 & 27.331 & 28.5 \\
Mar 2026 & 22.328 & 17.654 & 20.9 & 16.840 & 24.6 & 26.034 & 18.1 \\
\midrule
\textbf{Mean} & --- & --- & \textbf{16.2} & --- & \textbf{22.8} & --- & \textbf{29.9} \\
\bottomrule
\end{tabular}
\end{table}

\textbf{Best model: SARIMA} with a mean live error of \textbf{16.2\%}. LSTM achieves 22.8\%. ConvLSTM's errors are markedly higher (29.9\%), with January showing a 43\% deviation driven by the model over-predicting radiance in bright urban pixels. This highlights that the ConvLSTM model trained on the Kolkata data has not yet generalised the steep upward trend visible in the Kolkata 2026 actuals.

% ----- 6.3 ------------------------------------------------------------------
\vspace{0.5cm}
\begin{center}
    \hypertarget{sec:s63}{}%
    {\large \textbf{6.3 Discussion}}\\[0.3cm]
    \noindent\rule{\linewidth}{0.4pt}\\[0.5cm]
\end{center}

\textbf{Why SARIMA outperforms deep learning on monthly NTL.} The SARIMA$(0,1,1)(0,1,1)_{12}$ and SARIMA$(0,1,2)(0,1,1)_{12}$ orders selected for Kharagpur and Kolkata respectively are parsimonious models with well-matched seasonal structure. Monthly NTL time series are characterised by a dominant period-12 seasonality and weak non-linear dynamics. With only 120 training observations, the approximately 150,000-parameter LSTM and ConvLSTM models have insufficient data to reliably learn these patterns and generalise beyond the training distribution. SARIMA, with far fewer free parameters, is better regularised by the available data.

\textbf{Why Kolkata errors are higher.} Absolute radiance values in Kolkata are 2--3$\times$ higher than Kharagpur. The MAPE metric is therefore harder to minimise in absolute terms. Additionally, Kolkata's spatial heterogeneity (dense commercial core, residential periphery, industrial corridors) introduces more complex seasonal patterns than Kharagpur's campus-dominated signal.

\textbf{The unique value of ConvLSTM.} Despite its higher scalar error, the ConvLSTM provides information that SARIMA and LSTM cannot offer: \emph{pixel-level spatial forecasts}. The predicted frame identifies which sub-regions within the city are expected to brighten or dim, enabling spatially resolved urban monitoring. The SSIM values of 0.477 (Kharagpur) and 0.536 (Kolkata) indicate that the spatial structure of the predicted frames is meaningfully similar to the actual frames, even when the mean radiance value is off.

\textbf{Limitations.}
\begin{itemize}[leftmargin=1.5cm]
    \item Only two cities are included; results may not generalise to cities with very different NTL profiles.
    \item All models use only the NTL signal itself; no auxiliary predictors (temperature, rainfall, GDP, population) are included.
    \item The ConvLSTM is trained from scratch on a relatively small dataset (120 spatial sequences); pre-training on a larger global corpus would likely improve both accuracy and generalisation.
    \item The live validation covers only three months (January--March 2026); a longer evaluation period would yield more robust conclusions.
\end{itemize}

% =====================================================================
%  CHAPTER 7 — CONCLUSION AND FUTURE WORK
% =====================================================================
\newpage
\thispagestyle{plain}
\begin{center}
    \hypertarget{sec:ch7}{}%
    {\Large \textbf{7. Conclusion and Future Work}}\\[0.3cm]
    \noindent\rule{\linewidth}{0.6pt}\\[0.8cm]
\end{center}
\setlength{\parindent}{1cm}
\setlength{\parskip}{6pt}
\justifying

% ----- 7.1 ------------------------------------------------------------------
\vspace{0.5cm}
\begin{center}
    \hypertarget{sec:s71}{}%
    {\large \textbf{7.1 Conclusion}}\\[0.3cm]
    \noindent\rule{\linewidth}{0.4pt}\\[0.5cm]
\end{center}

This project has demonstrated the feasibility of multi-model nighttime light forecasting for Indian cities using freely available VIIRS satellite data from Google Earth Engine. Three models spanning statistical and deep learning approaches were implemented, trained, and rigorously evaluated on both a 24-month held-out test set and a three-month live validation against actual 2026 observations.

The key findings are as follows:

\begin{itemize}[leftmargin=1.5cm]
    \item \textbf{SARIMA is the best overall scalar forecaster} on both cities. It achieves a mean live validation error of \textbf{10.0\%} on Kharagpur and \textbf{16.2\%} on Kolkata, outperforming both deep learning alternatives on monthly NTL with 120 training observations. Its explicit seasonal parameterisation ($s=12$) aligns naturally with the dominant annual periodicity in the data.

    \item \textbf{LSTM provides a competitive deep-learning baseline} with live errors of 17.5\% (Kharagpur) and 22.8\% (Kolkata). While weaker than SARIMA on this dataset, it provides principled probabilistic predictions through MC Dropout uncertainty bands.

    \item \textbf{ConvLSTM adds a unique spatial forecasting dimension} not available from the scalar models. Despite higher mean-radiance error, its per-pixel spatial predictions carry structural information about which sub-regions of the city are changing --- information that is invisible to SARIMA and LSTM.

    \item All results are \textbf{validated against real 2026 data} downloaded independently in April 2026, confirming model reliability beyond the training distribution.

    \item The \textbf{Lumina dashboard}, deployed at \url{https://lumina-iamabhinav-dev.streamlit.app}, provides public, non-technical access to all underlying data, model outputs, and live validation results.
\end{itemize}

Beginning with the calibrated VIIRS reconstruction framework developed in BTP-1, this project has built a complete NTL analytics platform that extends from raw satellite imagery to interactive forecasting visualisations --- and lays the groundwork for the multi-modal research direction described in Section~7.3.

% ----- 7.2 ------------------------------------------------------------------
\vspace{0.5cm}
\begin{center}
    \hypertarget{sec:s72}{}%
    {\large \textbf{7.2 Future Work -- Near Term}}\\[0.3cm]
    \noindent\rule{\linewidth}{0.4pt}\\[0.5cm]
\end{center}

\begin{itemize}[leftmargin=1.5cm]
    \item \textbf{Extend to more Indian cities.} The pipeline is already city-parameterised via \texttt{src/cities.py}. Adding Mumbai, Delhi, Chennai, and Bengaluru requires only bounding box configuration and a data download step.

    \item \textbf{Longer live validation.} As more 2026 VIIRS composites become available (April--December 2026), the live validation window can be extended to provide a statistically more robust assessment.

    \item \textbf{Transformer-based temporal model.} The Temporal Fusion Transformer (TFT) architecture, designed explicitly for multi-horizon time series forecasting with covariate support, could be added as a fourth model to test whether attention-based mechanisms outperform LSTM on monthly NTL.

    \item \textbf{Improved ConvLSTM.} Incorporating attention gates in the decoder, or pre-training on a global VIIRS dataset before fine-tuning on city-specific data, would likely improve spatial forecasting accuracy on Kolkata.

    \item \textbf{GeoTIFF download in dashboard.} Currently the ConvLSTM forecast page provides GeoTIFF download only for historical frames. Enabling download of forecasted spatial grids in georeferenced format would increase the platform's utility for GIS practitioners.
\end{itemize}

% ----- 7.3 ------------------------------------------------------------------
\vspace{0.5cm}
\begin{center}
    \hypertarget{sec:s73}{}%
    {\large \textbf{7.3 Future Work -- Research Paper Proposal}}\\[0.3cm]
    \noindent\rule{\linewidth}{0.4pt}\\[0.5cm]
\end{center}

The most significant research direction enabled by this project is the integration of NTL time series with complementary socioeconomic data streams. The proposed work is framed as a future journal paper:

\begin{quote}
\textbf{Proposed title:}~\textit{``Multi-Modal Night-Time Light Analytics: Combining VIIRS, Population Projections, and GDP Estimates for Urban Growth Prediction in Indian Cities''}
\end{quote}

The key additions over the current BTP-2 work would be:

\begin{enumerate}[leftmargin=1.5cm]
    \item \textbf{Population data:} WorldPop gridded population datasets and census-based district-level population projections (2014--2026).
    \item \textbf{GDP proxy:} District-level Net State Domestic Product (NSDP) statistics and nighttime-light-based economic activity indices, providing an annual economic signal to complement the monthly NTL series.
    \item \textbf{Multi-modal architecture:} A fusion model that takes NTL time series, population growth rate, and GDP growth rate as input channels and forecasts future NTL (or directly forecasts GDP proxy) through a multi-input LSTM or Transformer.
    \item \textbf{Cross-city correlation study:} Analyse the lead-lag relationship between NTL growth, population growth, and GDP growth across a set of 10+ Indian cities to determine whether NTL leads or lags economic activity.
    \item \textbf{Policy application:} Use the trained multi-modal model to identify under-lit regions, forecast electricity demand, and detect economic anomalies in near-real-time.
\end{enumerate}

The Lumina platform provides precisely the infrastructure needed to pursue this direction. The calibrated NTL pipeline (via BTP-1 and BTP-2), the multi-city forecasting framework, and the deployed dashboard are all building blocks for a substantive multi-modal research contribution.

% =====================================================================
%  BIBLIOGRAPHY
% =====================================================================
\newpage
\thispagestyle{plain}
\begin{center}
    \hypertarget{sec:bib}{}%
    {\Large \textbf{Bibliography}}\\[0.3cm]
    \noindent\rule{\linewidth}{0.6pt}\\[0.8cm]
\end{center}

\bibliographystyle{unsrt}
\bibliography{references}

\end{document}
